%% kapa: exam=1 round=28 year=2017 period=1 form=A subject=economics subject_ko=경제학원론 seq=28 number=68
%% source: Q-Net 감정평가사 기출문제 (cst00309) artlSeq=5199603 file=extracted/1008/5199603/2017년 제28회 감정평가사 1차시험 문제지/1교시 A형.pdf
%% conversion: kordoc 4.24.51 -> tools/gen1.py (auto); review=proofread-sonnet; flags=none
\begin{question}{68}
신고전학파(Neoclassical) 투자이론에 관한 설명으로 옳지 않은 것은? (단, 모든 단위는 실질 단위이며 자본비용은 자본 한 단위당 비용이다.)
\begin{choices}
  \item 자본량이 증가하면 자본의 한계생산물은 감소한다.
  \item 감가상각률이 증가하면 자본비용도 증가한다.
  \item 자본량이 균제상태(steady state) 수준에 도달되면 자본의 한계생산물은 자본비용과 일치한다.
  \item 자본의 한계생산물이 자본비용보다 크다면 기업은 자본량을 증가시킨다.
  \item 실질이자율이 상승하면 자본비용은 감소한다.
\end{choices}
\end{question}
\answer{5}
